\documentclass[12pt,a4paper]{article}

\usepackage[utf8]{inputenc}
\usepackage[T2A]{fontenc}
\usepackage[russian]{babel}
\usepackage{geometry}
\usepackage{amsmath,amssymb}
\usepackage{graphicx}
\usepackage{hyperref}
\usepackage{booktabs}
\usepackage{array}
\usepackage{enumitem}
\usepackage{xcolor}
\usepackage{longtable}
\usepackage{float}

\geometry{margin=2.5cm}
\hypersetup{
    colorlinks=true,
    linkcolor=blue,
    urlcolor=blue,
    citecolor=blue
}

\title{GRA-Living-Vaccine-Architecture:\\
Концепция самообучающейся архитектуры живых терапий на основе ИИ}
\author{ }
\date{2026}

\begin{document}

\maketitle

\begin{abstract}
В работе предлагается концептуальная архитектура GRA-Living-Vaccine-Architecture (SLVA), предназначенная для описания самоадаптирующихся терапевтических систем, сочетающих методы искусственного интеллекта, синтетической биологии и иерархического управления стабильностью. В отличие от традиционных вакцин, рассматриваемая система трактуется как программируемый живой агент, способный реагировать на сигналы среды, исполнять логические правила и адаптировать поведение на основе обратной связи. Представлены уровни архитектуры, формальный DSL для задания поведения агента, прототип симуляционной среды и принципы безопасного nullification-механизма. Работа носит концептуальный характер и не содержит лабораторных протоколов.
\end{abstract}

\section{Введение}

Идея ``самообучающейся вакцины'' возникает на пересечении нескольких направлений: искусственного интеллекта, синтетической биологии, инженерии генетических цепей и адаптивных терапевтических систем. Современные исследования в области engineered living therapeutics и programmable biology демонстрируют, что клетки и биологические носители могут быть запрограммированы на выполнение терапевтических функций при наличии определённых внешних условий.

Однако большинство существующих работ фокусируется либо на биологической реализации, либо на моделировании отдельных компонентов. В данной статье предлагается более общая архитектурная рамка, в которой терапевтический агент рассматривается как иерархическая система с обратной связью, а ИИ выступает не как ``встроенный интеллект'' внутри препарата, а как внешний слой проектирования, оптимизации и обновления правил поведения.

\section{Постановка задачи}

Требуется описать архитектуру системы, которая:

\begin{enumerate}[label=\arabic*)]
    \item получает сигналы среды или организма;
    \item применяет набор логических правил;
    \item выполняет терапевтическое действие;
    \item контролирует собственную устойчивость;
    \item при нарушении ограничений переходит в безопасное состояние.
\end{enumerate}

Целью является построение абстрактной модели, пригодной для симуляции, прототипирования и последующего теоретического анализа, но не для прямой биомедицинской реализации.

\section{Общая архитектура}

Предлагаемая система состоит из пяти логических уровней.

\subsection{Уровень сигналов}

На этом уровне находятся входные данные:
\begin{itemize}
    \item уровень инфекции;
    \item воспаление;
    \item биомаркеры;
    \item внутренние состояния агента.
\end{itemize}

Сигналы могут быть как непрерывными, так и бинарными. Они формируют наблюдаемое состояние среды, которое далее обрабатывается логическим уровнем.

\subsection{Логический уровень}

Логический уровень задаёт правила вида:

\[
\text{IF condition} \rightarrow \text{action}
\]

Условия могут быть композиционными: например, несколько сигналов должны одновременно удовлетворять заданным порогам. Действия могут включать:
\begin{itemize}
    \item терапевтический эффект;
    \item изменение внутреннего состояния;
    \item эмиссию сигнала;
    \item nullification.
\end{itemize}

\subsection{Уровень состояния агента}

Агент обладает внутренними переменными:
\begin{itemize}
    \item энергия;
    \item активация;
    \item риск перегрузки;
    \item индикатор деактивации.
\end{itemize}

Состояние используется как часть условий правил, а также как основа для проверки устойчивости.

\subsection{Swarm-уровень}

Множество агентов могут взаимодействовать как распределённая группа. В этом случае система приобретает свойства swarm-архитектуры:
\begin{itemize}
    \item локальные решения;
    \item коллективная реакция;
    \item отсутствие единого центра исполнения;
    \item устойчивость к частичным сбоям.
\end{itemize}

\subsection{ИИ-ядро проектирования}

ИИ-ядро не является частью живого агента в прямом смысле. Оно служит внешним оптимизатором, который:
\begin{itemize}
    \item генерирует правила;
    \item подбирает параметры;
    \item оценивает устойчивость;
    \item обновляет архитектурные решения по результатам симуляции.
\end{itemize}

\section{Формальная модель}

Пусть агент задаётся кортежем:

\[
A = \langle S, R, X, \Pi \rangle
\]

где:
\begin{itemize}
    \item $S$ --- пространство состояний;
    \item $R$ --- набор правил;
    \item $X$ --- пространство сигналов;
    \item $\Pi$ --- политика устойчивости и деактивации.
\end{itemize}

Каждое правило можно записать как:

\[
r_i : c_i(X,S) \rightarrow a_i(S)
\]

где $c_i$ --- условие, а $a_i$ --- действие, изменяющее состояние агента.

Функция устойчивости может быть определена как:

\[
U(S) = \alpha \cdot E + \beta \cdot A + \gamma \cdot D
\]

где $E$ --- энергия, $A$ --- активация, $D$ --- мера опасности, а $\alpha, \beta, \gamma$ --- коэффициенты. Если $U(S)$ выходит за допустимый диапазон, применяется политика nullification.

\section{DSL для описания агента}

Для задания поведения агента предлагается простой предметно-ориентированный язык (DSL). Он позволяет описывать:
\begin{itemize}
    \item сигналы;
    \item состояние;
    \item правила;
    \item политику устойчивости.
\end{itemize}

Пример описания:

\begin{verbatim}
agent:
  id: toy-immunotherapy-agent
  signals:
    - infection_level
    - inflammation
  state:
    - energy
    - activation
  rules:
    - if infection_level > 0.7 and energy > 0.3
      then therapeutic_effect
  stability:
    max_activation_level: 1.0
    nullification_policy: hard
\end{verbatim}

Такой язык позволяет отделить архитектурное описание от механизма исполнения.

\section{Симуляционная среда}

Для проверки логики системы достаточно абстрактной симуляции. В ней:
\begin{enumerate}
    \item среда генерирует временные сигналы;
    \item агент читает сигналы;
    \item правила определяют действие;
    \item состояние обновляется;
    \item проводится проверка устойчивости.
\end{enumerate}

Цикл симуляции можно формализовать как:

\[
(X_t, S_t) \rightarrow a_t \rightarrow S_{t+1}
\]

где $X_t$ --- сигналы среды в момент времени $t$, $S_t$ --- состояние агента, $a_t$ --- действие.

\section{Механизм nullification}

Ключевая часть архитектуры --- безопасное обнуление активности при нарушении ограничений. Это необходимо для предотвращения неконтролируемой эскалации внутренней активации.

В простейшем случае nullification означает:
\begin{itemize}
    \item отключение правил;
    \item фиксацию состояния агента;
    \item переход в неактивный режим.
\end{itemize}

С точки зрения архитектуры это аналог защитного предохранителя.

\section{Обсуждение}

Предложенная модель полезна как концептуальный мост между:
\begin{itemize}
    \item ИИ-оптимизацией;
    \item синтетической биологией;
    \item формальными системами правил;
    \item архитектурами устойчивости.
\end{itemize}

Важно подчеркнуть, что речь идёт не о реальной вакцине с ИИ внутри, а о системе проектирования и симуляции программируемых живых терапий. Подобная рамка может использоваться для описания будущих медицинских агентов, но требует строгого разделения между теоретическим моделированием и биологической реализацией.

\section{Ограничения}

У данной работы есть ограничения:
\begin{itemize}
    \item отсутствие экспериментальной биологической валидации;
    \item упрощённость модели состояния;
    \item отсутствие стохастических иммунных процессов;
    \item отсутствие клинических параметров безопасности.
\end{itemize}

Поэтому данный материал следует рассматривать как исследовательский и архитектурный, а не как прикладное руководство.

\section{Заключение}

В работе предложена концептуальная архитектура GRA-Living-Vaccine-Architecture для описания самоадаптирующихся живых терапий. Архитектура объединяет уровни сигналов, логики, состояния, swarm-взаимодействия и ИИ-дизайна. Введён DSL для описания поведения агента и механизм безопасного nullification. Предложенная рамка может служить основой для дальнейших исследований в области программируемых терапий, симуляции биосистем и иерархических моделей устойчивости.

\section*{Благодарности}

Текст подготовлен как концептуальная база для дальнейшего развития архитектурного проекта.

\begin{thebibliography}{9}

\bibitem{syntheticbio}
Synthetic biology in medical and pharmaceutical applications: обзорные работы по программируемым биосистемам.

\bibitem{elt}
Engineered living therapeutics: исследования о живых терапевтических системах.

\bibitem{immunotherapy}
Synthetic biology approaches to immunotherapy: применение генетических цепей для управления терапевтическим поведением.

\bibitem{ai_design}
AI-assisted design of biological systems: использование ИИ для проектирования биологических последовательностей и схем.

\end{thebibliography}

\end{document}